\documentclass[11pt]{article}

\usepackage[T1]{fontenc}
\usepackage{amsmath,amssymb,amsthm}
\usepackage[margin=1in]{geometry}
\usepackage[hidelinks]{hyperref}

\newtheorem{theorem}{Theorem}
\newtheorem{lemma}[theorem]{Lemma}
\newtheorem{corollary}[theorem]{Corollary}
\newtheorem{conjecture}[theorem]{Conjecture}
\theoremstyle{remark}
\newtheorem{remark}[theorem]{Remark}

\newcommand{\R}{\mathbb R}
\newcommand{\D}{\mathcal D}
\newcommand{\WN}{\mathcal W^{\mathrm N}}
\newcommand{\WV}{\mathcal W^\nu}
\newcommand{\Cband}{\mathcal C}

\title{A Critical Common Budget for Type-II Frozen-Band Work}
\author{John Robert Lawson and Codex}
\date{Research round ending 24 July 2026}

\begin{document}
\maketitle

\begin{abstract}
This round answers the live summation question for the event-dependent
Gaussian tests on the energy-efficient Type-II branch.  The answer is
positive but scale-critical: nonlinear band work weighted by the square root
of the physical band radius, and squared viscous work weighted by radius
squared over event time, sum against the one physical dissipation measure.
After a scale-separated extraction, an exact Gaussian Bessel identity
extends the estimate to full first-correlation lifetimes.  The same
extraction exposes the ceiling of this approach: the fixed quantile gap
forces a uniform multiplicative band width, hence geometric radii and an
automatically summable nonlinear charge.  In the existing \(q_j=4\) power
ledger, viscous death is impossible infinitely often, but nonlinear rotation
remains compatible.  These are conditional analytic reductions, not a
regularity or breakdown theorem.
\end{abstract}

\section{Scope and setup}

Let \(u\) be a smooth finite-energy solution on
\([0,T^*)\times\R^3\):
\[
 \partial_tu+\mathbb P\nabla\!\cdot(u\otimes u)=\nu\Delta u,
 \qquad \nabla\!\cdot u=0.
\]
Set
\[
 E_*:=\sup_{t<T^*}\frac12\|u(t)\|_2^2,
 \qquad
 \D(I):=\nu\int_I\|\nabla u(t)\|_2^2\,dt,
 \qquad
 \D_*:=\D((0,T^*))\le E_*.
\]
The last inequality is the classical energy identity before \(T^*\).

At selected first-record times \(t_j\uparrow T^*\), retain the physical
Gaussian subgrid quantiles
\[
 \mathbf K(r_j^{\mathrm L},t_j)=\eta_{\mathrm L},
 \qquad
 \mathbf K(r_j^{\mathrm H},t_j)=\eta_{\mathrm H},
 \qquad
 \gamma:=\eta_{\mathrm H}-\eta_{\mathrm L}>0.
\]
For \(G_rf=\Gamma_r*f\), define
\[
 S_j:=G_{r_j^{\mathrm L}}^2-G_{r_j^{\mathrm H}}^2,
 \qquad
 B_j:=S_j^{1/2},
 \qquad
 \psi_j:=S_ju(t_j).
\]
The multiplier
\[
 \mathfrak s_j(\xi)
 =
 e^{-(r_j^{\mathrm L})^2|\xi|^2}
 -
 e^{-(r_j^{\mathrm H})^2|\xi|^2}
\]
is nonnegative and
\[
 \int_{\R^3}\mathfrak s_j|\widehat u(t_j)|^2\,d\xi
 =
 \|B_ju(t_j)\|_2^2
 =
 2\gamma.
\]

For \(s\in(t_j,T^*)\), put
\[
\begin{aligned}
 \WN_{j,s}
 &:=
 \int_{t_j}^{s}\!\!\int_{\R^3}
 u\otimes u:\nabla\psi_j\,dx\,dt,\\
 \WV_{j,s}
 &:=
 -\nu\int_{t_j}^{s}\!\!\int_{\R^3}
 \nabla u:\nabla\psi_j\,dx\,dt.
\end{aligned}
\]
Then
\[
 \Cband_j(s)-2\gamma
 =
 \WN_{j,s}+\WV_{j,s},
 \qquad
 \Cband_j(s):=
 \langle B_ju(s),B_ju(t_j)\rangle.
\]
All conclusions below are conditional on the branch that produced these
fixed-gap quantiles.  No argument derives that branch from arbitrary Clay
data.

\section{Two critical estimates}

\begin{lemma}[Fractional norm of the frozen band]
For every \(\alpha\ge0\),
\[
 \|\psi_j\|_{\dot H^\alpha}^2
 \le
 C_\alpha\gamma(r_j^{\mathrm L})^{-2\alpha}.
\]
In particular,
\[
 \|\psi_j\|_{\dot H^{1/2}}
 \lesssim
 \sqrt{\frac{\gamma}{r_j^{\mathrm L}}},
 \qquad
 \|\nabla\psi_j\|_2
 \lesssim
 \frac{\sqrt\gamma}{r_j^{\mathrm L}}.
\]
\end{lemma}

\begin{proof}
Write
\[
 x=(r_j^{\mathrm L})^2|\xi|^2,
 \qquad
 \rho_j=r_j^{\mathrm H}/r_j^{\mathrm L}>1.
\]
Then
\[
 0\le\mathfrak s_j=e^{-x}-e^{-\rho_j^2x}\le e^{-x}.
\]
Since \(x^\alpha e^{-x}\le C_\alpha\),
\[
 x^\alpha\mathfrak s_j^2
 \le C_\alpha\mathfrak s_j.
\]
Consequently
\[
\begin{aligned}
 \|\psi_j\|_{\dot H^\alpha}^2
 &=
 \int |\xi|^{2\alpha}
 \mathfrak s_j^2|\widehat u(t_j)|^2\,d\xi\\
 &\le
 C_\alpha(r_j^{\mathrm L})^{-2\alpha}
 \int\mathfrak s_j|\widehat u(t_j)|^2\,d\xi\\
 &=
 2C_\alpha\gamma(r_j^{\mathrm L})^{-2\alpha}.
\end{aligned}
\]
\end{proof}

\begin{lemma}[Critical nonlinear dual estimate]
For every solenoidal \(v\in H^1(\R^3)\),
\[
 \|\mathbb P\nabla\!\cdot(v\otimes v)\|_{\dot H^{-1/2}}
 \lesssim
 \|\nabla v\|_2^2.
\]
Hence, for solenoidal \(\phi\in\dot H^{1/2}\),
\[
 \left|\int_{\R^3}v\otimes v:\nabla\phi\,dx\right|
 \lesssim
 \|\nabla v\|_2^2\|\phi\|_{\dot H^{1/2}}.
\]
\end{lemma}

\begin{proof}
The Leray projector is a zero-order multiplier and
\(\dot H^{-1/2}\nabla\) has order \(1/2\).  Thus
\[
 \|\mathbb P\nabla\!\cdot(v\otimes v)\|_{\dot H^{-1/2}}
 \lesssim
 \|v\otimes v\|_{\dot H^{1/2}}.
\]
The fractional Leibniz inequality, with spatial exponents \(3\) and \(6\),
followed by Sobolev embedding, gives
\[
\begin{aligned}
 \|v\otimes v\|_{\dot H^{1/2}}
 &\lesssim
 \|D^{1/2}v\|_3\|v\|_6\\
 &\lesssim
 \|\nabla v\|_2^2.
\end{aligned}
\]
Duality proves the final assertion.
\end{proof}

\section{The common weighted budget}

Put
\[
 I_j=[t_j,t_{j+1}],
 \qquad
 h_j=t_{j+1}-t_j,
 \qquad
 r_j=r_j^{\mathrm L},
\]
and abbreviate
\(\WN_j=\WN_{j,t_{j+1}}\),
\(\WV_j=\WV_{j,t_{j+1}}\).

\begin{theorem}[Common consecutive-record budget]
\[
\boxed{
 \sum_j\left[
 \nu\sqrt{\frac{r_j}{\gamma}}\,|\WN_j|
 +
 \frac{r_j^2|\WV_j|^2}{\gamma\nu h_j}
 \right]
 \lesssim
 \D_*.
}
\]
\end{theorem}

\begin{proof}
The preceding lemmas yield
\[
\begin{aligned}
 |\WN_j|
 &\lesssim
 \|\psi_j\|_{\dot H^{1/2}}
 \int_{I_j}\|\nabla u(t)\|_2^2\,dt\\
 &\lesssim
 \sqrt{\frac{\gamma}{r_j}}\,
 \frac{\D(I_j)}{\nu}.
\end{aligned}
\]
Therefore
\[
 \nu\sqrt{\frac{r_j}{\gamma}}\,|\WN_j|
 \lesssim \D(I_j).
\]
For viscosity, Cauchy--Schwarz in time gives
\[
\begin{aligned}
 |\WV_j|
 &\le
 \nu\|\nabla\psi_j\|_2
 h_j^{1/2}
 \left(\int_{I_j}\|\nabla u(t)\|_2^2\,dt\right)^{1/2}\\
 &\lesssim
 \frac{\sqrt\gamma}{r_j}
 \sqrt{\nu h_j\D(I_j)}.
\end{aligned}
\]
Hence
\[
 \frac{r_j^2|\WV_j|^2}{\gamma\nu h_j}
 \lesssim\D(I_j).
\]
The intervals \(I_j\) are pairwise disjoint, so their dissipation measures
sum to at most \(\D_*\).
\end{proof}

\begin{corollary}[Additive correlation-loss charges]
Let
\[
 \mathcal L=\{j:\Cband_j(t_{j+1})<\gamma\}.
\]
There is a disjoint partition
\(\mathcal L=\mathcal N\mathbin{\dot\cup}\mathcal V\) such that
\[
 |\WN_j|\ge\gamma/2\quad(j\in\mathcal N),
 \qquad
 |\WV_j|\ge\gamma/2\quad(j\in\mathcal V),
\]
and
\[
\boxed{
 \sum_{j\in\mathcal N}\nu\sqrt{\gamma r_j}
 +
 \sum_{j\in\mathcal V}
 \frac{\gamma r_j^2}{\nu h_j}
 \lesssim
 \D_*.
}
\]
In particular,
\[
 \sum_{j\in\mathcal L}
 \min\left\{
 \nu\sqrt{\gamma r_j},
 \frac{\gamma r_j^2}{\nu h_j}
 \right\}
 <\infty.
\]
\end{corollary}

\begin{proof}
For \(j\in\mathcal L\), the exact correlation identity gives
\[
 |\WN_j+\WV_j|>\gamma.
\]
Assign the event to a term of magnitude at least \(\gamma/2\), then invoke
the theorem.
\end{proof}

\begin{remark}
The theorem genuinely uses one physical quantity, \(\D_*\), for every
event-dependent test.  It does not bound \(\sum_j|\WN_j|\): the critical
dual norm necessarily contributes \(r_j^{-1/2}\).
\end{remark}

\section{First-death lifetimes}

Because \(r_j^{\mathrm H}\to0\), choose an infinite subsequence, relabelled
by \(k\), for which
\[
 r_{k+1}^{\mathrm H}<r_k^{\mathrm L}.
\]

\begin{lemma}[Exact Gaussian Bessel packing]
For every \(f\in L^2(\R^3)\),
\[
 \sum_k\|B_kf\|_2^2\le\|f\|_2^2.
\]
\end{lemma}

\begin{proof}
The multiplier has the heat-scale representation
\[
 \mathfrak s_k(\xi)
 =
 \int_{(r_k^{\mathrm L})^2}^{(r_k^{\mathrm H})^2}
 |\xi|^2e^{-a|\xi|^2}\,da.
\]
The integration intervals are disjoint, so
\[
 \sum_k\mathfrak s_k(\xi)
 \le
 \int_0^\infty|\xi|^2e^{-a|\xi|^2}\,da
 =1.
\]
Apply Plancherel.
\end{proof}

Define the first death time
\[
 \sigma_k
 :=
 \inf\{t\in(t_k,T^*):\Cband_k(t)\le\gamma\},
\]
with \(\sigma_k=T^*\) if the set is empty.  The band is alive on
\([t_k,\sigma_k)\).

\begin{theorem}[Bounded-overlap lifetime budget]
At each \(t<T^*\), at most
\[
 M:=\left\lceil\frac{4E_*}{\gamma}\right\rceil
\]
separated bands are alive.  At most \(M\) never die.  The finite lifetime
intervals \(J_k=[t_k,\sigma_k]\) have overlap at most \(M\), and
\[
\boxed{
 \sum_k\left[
 \nu\sqrt{\frac{r_k^{\mathrm L}}{\gamma}}\,
 |\WN_{k,\sigma_k}|
 +
 \frac{
 (r_k^{\mathrm L})^2|\WV_{k,\sigma_k}|^2
 }{
 \gamma\nu(\sigma_k-t_k)
 }
 \right]
 \lesssim
 M\D_*.
}
\]
After assigning every finite death to a work floor,
\[
\boxed{
 \sum_{k\in\mathcal N_{\rm death}}
 \nu\sqrt{\gamma r_k^{\mathrm L}}
 +
 \sum_{k\in\mathcal V_{\rm death}}
 \frac{
 \gamma(r_k^{\mathrm L})^2
 }{
 \nu(\sigma_k-t_k)
 }
 \lesssim
 M\D_*.
}
\]
\end{theorem}

\begin{proof}
If \(k\) is alive at \(t\), Cauchy--Schwarz and
\(\|B_ku(t_k)\|_2^2=2\gamma\) imply
\[
 \|B_ku(t)\|_2^2>\frac\gamma2.
\]
The Bessel lemma and \(\|u(t)\|_2^2\le2E_*\) therefore give
\[
 \frac\gamma2\,\#\{k:k\text{ alive at }t\}
 <
 \sum_{k:\,\mathrm{alive}}\|B_ku(t)\|_2^2
 \le2E_*.
\]
This proves the count and the bound on immortal bands.  It also bounds the
overlap of the finite \(J_k\).

At finite death, continuity gives
\(\Cband_k(\sigma_k)=\gamma\), so one work has magnitude at least
\(\gamma/2\).  The proof of the consecutive-record theorem applies on
each \(J_k\).  Finally,
\[
 \sum_k\D(J_k)\le M\D_*.
\]
\end{proof}

\section{The square-function ceiling}

\begin{lemma}[Uniform multiplicative band width]
Every fixed-gap quantile band obeys
\[
 \frac{r_j^{\mathrm H}}{r_j^{\mathrm L}}
 \ge
 \Lambda
 :=
 \left(1+\frac{e\gamma}{E_*}\right)^{1/2}>1.
\]
\end{lemma}

\begin{proof}
Put \(\rho=r_j^{\mathrm H}/r_j^{\mathrm L}\).  For \(x\ge0\),
\[
 e^{-x}-e^{-\rho^2x}
 =
 \int_1^{\rho^2}xe^{-sx}\,ds
 \le
 (\rho^2-1)xe^{-x}
 \le
 \frac{\rho^2-1}{e}.
\]
It follows that
\[
 2\gamma
 \le
 \frac{\rho^2-1}{e}\|u(t_j)\|_2^2
 \le
 \frac{2E_*(\rho^2-1)}{e}.
\]
\end{proof}

For the separated subsequence,
\[
 r_{k+1}^{\mathrm H}
 <
 r_k^{\mathrm L}
 \le
 \Lambda^{-1}r_k^{\mathrm H}.
\]
Thus
\[
 \sum_k\sqrt{r_k^{\mathrm L}}<\infty.
\]
The extraction required for exact Bessel packing therefore makes the
critical nonlinear death charge automatically summable.  Bessel packing,
the energy identity, and critical fractional duality alone cannot exclude
infinitely many nonlinear deaths.

\section{Audit of the exact power survivor}

The current representative has
\[
 r_j^{\mathrm L}\asymp2^{-5j},
 \qquad
 h_j\asymp\frac{2^{-11j}}j,
 \qquad
 \nu=1.
\]
Its nonlinear and viscous death charges are respectively
\[
 \nu\sqrt{\gamma r_j^{\mathrm L}}
 \asymp
 \sqrt\gamma\,2^{-5j/2},
\]
and
\[
 \frac{\gamma(r_j^{\mathrm L})^2}{\nu h_j}
 \asymp
 \gamma j2^j.
\]
Therefore
\[
 \sum_j\sqrt{r_j^{\mathrm L}}<\infty,
 \qquad
 \sum_j\frac{(r_j^{\mathrm L})^2}{h_j}=\infty.
\]
All but finitely many correlation deaths in that ledger must be nonlinear.
The old abstract Hilbert rotation remains compatible with this assignment:
its residence reservoir \(c_j\asymp2^{-j}/j\) dominates the new
\(2^{-5j/2}\) charge.  The rotation is still not a Navier--Stokes velocity.

\section{Robust findings, subject to review}

\begin{enumerate}
\item Frozen Gaussian tests have exact scale-adapted fractional Sobolev
ceilings determined by their fixed energy gap.
\item Critical \(\dot H^{-1/2}\)--\(\dot H^{1/2}\) duality converts all
changing nonlinear tests into one weighted dissipation budget.
\item Viscous works possess a simultaneous quadratic weighted budget.
\item Gaussian heat-scale separation gives an exact Bessel inequality and a
finite bound on simultaneous correlation lifetimes.
\item The fixed gap forces uniform logarithmic band width, so elementary
spectral thinning cannot make the nonlinear charge nonsummable.
\item The \(q_j=4\) survivor is forced into the nonlinear-death branch.
\end{enumerate}

The constituent inequalities are standard or proved above.  Their assembly
for this Type-II record stack has not received external mathematical review;
no literature-novelty claim is made here.

\section{Conjectures and open obligations}

\begin{conjecture}[Unthinned heat-scale Carleson law]
The complete, overlapping family of first-record heat-scale intervals obeys
a same-trajectory Carleson or variation estimate that retains event density
without the geometric thinning used by the Bessel theorem.
\end{conjecture}

\begin{conjecture}[Triadic rotation obstruction]
The Navier--Stokes quadratic interaction cannot realise infinitely many
fixed-gap nonlinear correlation deaths while paying only the summable
\(\sqrt{r_j}\) dissipation charges allowed by critical fractional duality.
\end{conjecture}

\begin{conjecture}[Critical sharpness]
Without additional trajectory, spatial, or sign structure, the
half-derivative loss and the resulting \(\sqrt{r_j}\) event charge cannot be
improved by energy-class estimates.
\end{conjecture}

The next obligations are:
\begin{enumerate}
\item analyse the unthinned overlap multiplicity as a measure on logarithmic
heat scale;
\item seek a temporal variation or adjoint cancellation beyond the critical
fractional product estimate;
\item otherwise build or exclude an actual NSE triadic analogue of the
nonlinear rotation;
\item separately retain spatial recentering, coherent-trace rigidity, and
the divergent-normalised-energy cells.
\end{enumerate}

\section{Epistemic conclusion}

This round closes the existence of a common \emph{weighted} band-work
budget and closes the elementary scale-separated Bessel route negatively.
It does not give an unweighted work sum, rule out the remaining nonlinear
rotation, establish regularity or breakdown, or prove any Clay alternative
A--D.

The proof-owning note is stored as
\begin{center}
\small
\texttt{dossier/experiments/}\\
\texttt{type-ii-band-dissipation-budget.md}.
\end{center}
The note was retired in the 2026-07-27 curation; it is recoverable from
Git history before commit \texttt{13557eb}, per the recipe in
\texttt{dossier/records/README.md}.
Canonical route and experiment metadata remain in
\texttt{dossier/records/}.  The public repository is
\url{https://github.com/Bingham-Research-Center/brc-navier-stokes}.

\end{document}
